\documentclass[12pt,a4paper]{article}

\usepackage[english]{babel}
\usepackage[utf8]{inputenc}
\usepackage{amsmath}
\usepackage{amssymb}
\usepackage{graphicx}
\usepackage{hyperref}

\begin{document}
\thispagestyle{empty}
\begin{center}
Rheinisch-Westf\"alische Technische Hochschule Aachen \\
Alessandro Berti\\[6ex]

\LARGE
	\textbf{Project Initialization: Implementation of "Conformance Predictions"} \\[6ex]
        \centering{Group 4:}\\
	\textit{Richard Schuppener, Cem Colak} \\[6ex]
\end{center}

\newpage

\tableofcontents
\newpage

\section{Introduction}

\section{Business Case} 

\subsection{Client}
\subsection{Example Use Case 1:}
\subsection{Example Use Case 2:}

\section{Feasibility Study}

\subsection{Theoretical Feasibility}
\subsection{Technical Feasibility}

\section{Project Charter}\label{sec:charter} 

\includegraphics[width=1\textwidth]{project_charter.png}
\\\\
The milestones and their respective deadlines for submission of the required materials are:
\begin{itemize}
	\item Project Setup:
	\item Requirements Engineering: 
	\item Sprint 1: 
	\item Sprint 2: 
	\item Sprint 3: 
	\item Testing and Assessment: 
\end{itemize}
It is essential to acknowledge that the schedule presented here is preliminary and may be subject to adjustments. This flexibility accommodates unforeseen personal or professional circumstances that may arise during the project's execution, allowing team members to navigate challenges effectively. 

\section{Tools}
\subsection{Trello}
In order to plan, organize, and track the development of our work, we rely on a project management tool called Trello. It is based on the concept of boards, lists, and cards, allowing us to break down complex tasks into smaller ones and distribute them among team members. We utilize features such as due dates and labels to further customize our tasks and distinguish them from one another. This web-based platform makes it easy for group members to monitor the progress of each task and identify their responsibilities. In conclusion, Trello's user-friendly interface and comprehensive features significantly contribute to the success of our project.
\\\\
Link to our Trello board: \href{https://trello.com/b/Hlx8Wp8m}{\textit{Group 4 Trello}}

\subsection{Git and GitLab}
Version control systems are indispensable for modern software development, especially in collaborative environments like academic projects. For our PM4PY based project we used Git and GitLab, which will be introduced in the following: \\ 
\\
\textit{Git}, at its core, is a distributed version control system that tracks changes in source code over time. Its key strengths lie in its ability to create branches for remote independent development, merge code easily, and revert to previous versions. This allows team members to work concurrently on different features or bug fixes without interfering with each other's progress. Remote repositories enable the developing team to get the most up-to-date version of the project and develop from there. \\
\\
\textit{GitLab} on the other hand, provides a web-based platform built around Git, extending its functionalities with a suite of project management and collaboration tools. These include issue tracking, code review, continuous integration/continuous delivery (CI/CD) pipelines, and wiki pages for documentation. By integrating these features, GitLab creates a centralized hub for project activity.
\\\\

\subsection{PM4Py}
PM4Py is an open source library designed by the process mining group of the Frauenhofer Institute (FiT), which includes a plethora of state-of-the-art process mining algorithms. In regards to our project and the overall topic of Conformance Checking, PM4Py supports f.e. token-based, alignment-based and footprint-based conformance diagnostics, while also offering a lot of visualization techniques. Filtering event logs and extracting features is also supported, which makes pm4py it perfect for our project.


\section{Team}
\subsection{Personal Description}
\subsubsection{Richard Schuppener}
\subsubsection{Cem Colak}

\section{Roles}

\underline{Scrum Master}: 
\\\\
\underline{Official Communication/Product Owner}: 
\\\\
\underline{Tools Administration}: 
\\\\
\underline{Theory Expert}: 
\\\\
\underline{Software Expert}: 
\\\\
\underline{Test Management}: 


\section{Phase review}
\subsection{Richard Schuppener}
\subsection{Cem Colak}

\paragraph{}
\begin{thebibliography}{10}

\bibitem{denisov2018performance}
Denisov, V., Belkina, E., Fahland, D., van der Aalst, W. M. (2018, September). The Performance Spectrum Miner: Visual Analytics for Fine-Grained Performance Analysis of Processes. In BPM (Dissertation/Demos/Industry) (pp. 96-100).

\end{thebibliography}


\end{document}
